\begin{figure}[tb]
\centering
\begin{tikzpicture}[x=1.1cm,y=1.1cm,>=Latex,font=\small]
 \fill[blue!6] (-3,.75)--(0,.75)--(0,2.15)--(3,2.15)--(3,3.55)--(-3,3.55)--cycle;
 \draw[->] (-3.45,0)--(3.55,0) node[right] {$x$};
 \draw[->] (-3.45,0)--(-3.45,3.85) node[above] {$\log s$};
 \draw[blue!65!black,very thick] (-3,.75)--(0,.75)--(0,2.15)--(3,2.15);
 \draw[densely dashed,gray] (-3.35,3.3)--(3.1,3.3);
 \node[left] at (-3.45,3.3) {$\log\rho$};
 \node[blue!45!black] at (1.5,2.7) {continuation region};
 \node[anchor=east,align=right] at (-.5,1.9) {vertical side of the\\entrance boundary};
 \draw[->,gray] (-.48,1.7)--(-.08,1.6);
 \draw[orange!75!black,thick,->] (.05,1.3)--(.05,3.13);
 \fill (0,1.27) circle (2pt);
 \node[right] at (.22,1.22) {contact at $\underline s$};
 \draw[fill=white,blue!65!black,thick] (0,2.15) circle (2pt);
 \node[right,fill=white,inner sep=2pt] at (.25,2.02) {threshold equality at $s_0$};
\end{tikzpicture}
\caption{A schematic entrance boundary. In this illustration the contact
scale is denoted by $\underline s$ and a larger threshold scale by $s_0$.
A contact on a vertical side can have
derivative scale $\underline s<s_0$, where $s_0$ is a larger threshold scale.
The contact rule controls the initial derivative and the cumulative integral
along the outward vertical. It does not impose a pointwise sign at every later
scale. The two scales remain distinct in the final mass charge.}
\label{fig:contact}
\end{figure}
